%% Theory of Computing
%% File: v016a002.tex
%% Publication date: September 7, 2020
%% Authors: Andrej Bogdanov, Siyao Guo, and Ilan Komargodski

%

%
%
%
%

%

%
%
%
%
%
%
%
%
%
%
%
%
%
%
%
%
%
%
%
%
%
%
%
%
%
%
%
%
%

%

\documentclass{toc}

\DeclareMathOperator\pr{\mathop{\bf Pr}}
\DeclareMathOperator\E{\mathop{\bf E}}
\DeclareMathOperator\D{\mathrm{D}}
\DeclareMathOperator\Var{\mathrm{Var}}
\DeclareMathOperator\bias{\mathrm{Bias}}
\newcommand\unorm{\mathit{U}}


%
%

%
%
%
\tocdetails{%
  title = {Threshold Secret Sharing Requires a Linear-Size Alphabet},
%
%
%
%
  number_of_pages = {17}, 
  number_of_bibitems = {37},
  number_of_figures = {0},
  conference_version = {TCC16}, %
    %
    %
%
%
%
  author = {Andrej Bogdanov, Siyao Guo, and Ilan Komargodski},
    %
    %
    %
  authorlist = {Andrej Bogdanov, Siyao Guo, Ilan Komargodski},
    %
    %
    %
    %
    %
    %
    %
%
%
  %
  %
  %
  %
%
  acmclassification = {F.1.3},
    %
    %
  amsclassification = {68Q17, 94A60},
    %
    %
  keywords = {secret sharing, threshold, lower bound},
    %
    %
    %
    %
    %
}   %

    %
    %
    %
    %
    %
    %
    %
    %
    %
    %
    %
    %
    %
    %
    %
    %
    %
    %
    %
    %
    %
    %
    %
    %
    %
    %
    %
    %
    %

%
%
\tocdetails{%
volume = 16,
number = 2,
year = 2020,
%
%
%
%
%
%
%
received = {November 26, 2016},
revised = {May 3, 2019},
published = {September 7, 2020},
%
%
doi = {10.4086/toc.2020.v016a002},  
%
}   %

\begin{document}

\begin{frontmatter}[classification=text]
  
%
\title{Threshold Secret Sharing \\ Requires a Linear-Size Alphabet}  %

%
%

\title{Threshold Secret Sharing \\ Requires a Linear-Size Alphabet\titlefootnote{A conference version of this paper appeared in
the \href{https://link.springer.com/chapter/10.1007/978-3-662-53644-5_18}{Proceedings of the The 14th Theory of Cryptography Conference (TCC), 2016-B}.}}

%
%
%
%
%
%
%
%
%
%
%
%


%
%
%
%

%

\author[Bogdanov]{Andrej Bogdanov\thanks{Supported by RGC GRF grants  CUHK410113 and CUHK14208215.}}
\author[Guo]{Siyao Guo\thanks{Part of the work done in the Chinese University of Hong Kong supported by RGC GRF grants CUHK410112 and CUHK410113.}}
\author[Komargodski]{Ilan Komargodski\thanks{Supported by an AFOSR grant
FA9550-15-1-0262. Part of this work done while
    visiting CUHK, supported by RGC GRF grant CUHK410113, and while being a Ph.D.\
    student at the Weizmann Institute of Science, supported in part by a
    Levzion fellowship, by a grant from the I-CORE Program of the Planning and
    Budgeting Committee, the Israel Science Foundation, BSF and the Israeli
    Ministry of Science and Technology.}}

%
%
%
%

%
%
%
%
%
%
\begin{abstract}
  We prove that for every $n$ and $1 < t < n$ any $t$-out-of-$n$ threshold
  secret sharing scheme for one-bit secrets requires share size $\log(t + 1)$.
  Our bound is tight when $t = n - 1$ and $n$ is a prime power.  In 1990 Kilian
  and Nisan proved the incomparable bound $\log(n - t + 2)$.  Taken together,
  the two bounds imply that the share size of Shamir's secret sharing scheme
  (Comm.\ ACM 1979) is optimal up to an additive constant even for one-bit
  secrets for the whole range of parameters $1 < t < n$.
	
  More generally, we show that for all $1 < s < r < n$, any ramp secret sharing
  scheme with secrecy threshold $s$ and reconstruction threshold $r$ requires
  share size $\log((r + 1)/(r - s))$.
	
  As part of our analysis we formulate a simple game-theoretic relaxation of
  secret sharing for arbitrary access structures.  We prove the optimality of
  our analysis for threshold secret sharing with respect to this method and
  point out a general limitation.
\end{abstract}

\iffalse %
%
%
%
%
%
%
%
%
%
%
%
%
%
%
%
%
%
%
%
%
%

\tocarxivcategory{cs.CC,cs.CR}

\fi %

\end{frontmatter}

%
%
%

\newcommand{\party}{\mathsf p}
\newcommand{\parties}{\mathcal P}
\newcommand{\share}{\sigma}
\newcommand{\shares}{\sigma}
\newcommand{\SHARE}{\mathsf{SHARE}}
\newcommand{\RECON}{\mathsf{RECON}}

\newcommand{\secret}{s}
\newcommand{\thr}[2]{\mathsf{THR}_{#1}^{#2}}
\newcommand{\ramp}[2]{\mathsf{RAMP}_{#1}^{#2}}
\newcommand\F{\mathbb F}
\def\B{\{0,1\}}
\def\eps{\varepsilon}
\def\xor{\oplus}
\def\implies{\Rightarrow}
\def\cond{\mid}
\def\sign{\mathrm{sign}}

\def\calA{\mathcal{A}}
\def\calB{\mathcal{B}}
\def\calS{\mathcal{S}}
\def\calSmax{\hat{\mathcal{S}}}
\def\calR{\mathcal{R}}
\def\calP{\mathcal{P}}

\providecommand\floor[1]{\lfloor#1\rfloor}
\providecommand\ceil[1]{\lceil#1\rceil}
\providecommand\blog[1]{\log_2\ceil{#1}}
\providecommand\abs[1]{\lvert#1\rvert}
\providecommand\compl[1]{#1^{\complement}}
\providecommand\norm[1]{\|#1\|}
\providecommand\bigabs[1]{\bigl\lvert#1\bigr\rvert}
\providecommand\ip[1]{\langle#1\rangle}



\section{Introduction}
In 1979, Shamir~\cite{Shamir79} and Blakley~\cite{Blakley79} presented a method
for sharing a piece of secret information among $n$ parties such that any
$1<t<n$ parties can recover the secret while any $t-1$ parties learn 
	nothing about the secret. These methods are called $(t,n)$-\emph{threshold
	secret sharing schemes}. This sharp threshold between secrecy and
reconstruction is fundamental in applications where a group of mutually
suspicious individuals with conflicting interests must cooperate. Indeed,
threshold secret sharing schemes have found many applications in cryptography
and distributed computing; see the extensive survey of Beimel~\cite{Beimel11}
and the recent book of Cramer et al.~\cite{CDN2015}.

Threshold secret sharing was generalized by Ito et al.~\cite{ISN93} to 
allow more general structures of subsets to learn the
secret, while keeping the secret perfectly hidden from all other subsets. The
collection of qualified subsets is called an \emph{access structure}.

A significant goal in secret sharing is to minimize the share size, namely, the
amount of information distributed to the parties. Despite the long history of the
subject, there are significant gaps between lower and upper bounds both
for general access structures and for the special case of threshold structures.

\paragraph{Threshold access structures.} For $(t,n)$-threshold access
structures (denoted by $\thr{t}{n}$) and a 1-bit secret, Shamir~\cite{Shamir79} gave a very elegant
and efficient scheme: the dealer picks a random polynomial of degree $t-1$
conditioned on setting the free coefficient to be the secret, and gives the
$i$-th party the evaluation of the polynomial at the point $i$. The computation
is done over a field $\F$ of size $q>n$.

The correctness follows because one can recover the unique polynomial from any
$t$ points (and thus recover the secret). Security follows by a counting
argument showing that given less than $t$ points, all possibilities for the free
coefficient are equally likely. The share of each party is an element in the
field $\F$ that can be represented using $\log q\approx \log n$ bits. (All our
logarithms are in base 2.) The efficiency of this scheme makes it very attractive
for applications.

A natural question to ask is whether $\log{n}$-bit shares are necessary
for sharing a 1-bit secret for threshold access structures. Kilian and
Nisan~\cite{KilianN90}\footnote{Their result is unpublished and independently
	obtained (and generalized in various ways) by~\cite{PueyoCX13}. The original
	argument of Kilian an Nisan appears in \cite[Appendix A]{PueyoCX13} and was
	referenced earlier in~\cite{BeimelC94,BeimelPHD,BeimelF07}.}  showed that
$\log n$ bits are necessary when $t$ is not too large. Specifically, they proved
a $\log(n-t+2)$ lower bound on share size for $(t,n)$-threshold
schemes.  For large values of $t$, especially those close to $n$, their
bound does not rule out schemes with shares much shorter than $\log n$ bits. Their
bound leaves open the possibility that, in particular, $(n-1,n)$-threshold schemes
with two-bit shares exist.  


\emph{Ramp schemes} are a generalization of threshold schemes that allow for a 
gap between the secrecy and reconstruction parameters. In an $(s,r,n)$-ramp scheme, we
require that any subset of at least $r$ parties can recover the secret, while
any subset of size at most $s$ cannot learn anything about the
secret.\footnote{Another common definition (see~\cite[Definition 2.7]{FarrasHKP14} and~\cite[Example 2.11]{MolleviMP16} for examples) for a ramp scheme is where the
	information about the secret increases with the size of the set. 
We focus only
on the definition in which sets of size below a certain threshold have no
information about the secret, while sets of size larger than some threshold
can recover it.}  When $r=s+1$, an $(s,r,n)$-ramp scheme is exactly an
$(r,n)$-threshold scheme. Ramp schemes, defined by Blakley and
Meadows~\cite{BlakleyM84}, are useful for various applications (see,
\eg, \cite{FranklinY92,StinsonW99,ChenC06,MartinPS11}). If $r-s$ is large, they
can sometimes be realized with shorter shares than standard threshold schemes
(especially in the case of a long secret).

Generalizing the lower bound of Kilian and Nisan, Cascudo et al.~\cite{PueyoCX13} 
showed that $\log((n-s+1)/(r-s))$-bit shares are necessary to realize an
$(s,r,n)$-ramp scheme.   When $s = n - O(1)$, however, their share size bound
is a constant independent of $n$. Paterson and Stinson~\cite{PatersonS13} showed that this bound is tight for specific small values of $s$.

\paragraph{General access structures.}
For most access structures, the best known secret sharing schemes require shares
of size $2^{O(n)}$ for sharing a 1-bit secret. Specifically, viewing the access
structure as a Boolean indicator function for qualified subsets, the schemes of~\cite{ISN93,Leichter88,KarchmerW93} result with shares of size proportional to
the DNF/CNF size, monotone formula size, or monotone span program size of the
function, respectively. Thus, even for many access structures that can be
described by a small monotone uniform circuit, the best schemes have super-polynomial size shares.\footnote{One example is the \emph{directed
connectivity} access structure: the parties correspond to edge slots in the
complete \emph{directed} graph and the qualified subsets are those edges that
connect two distinguished nodes $s$ and $t$.}  On the other hand, the best
known lower bound on share size for sharing an $\ell$-bit secret is 
$\Omega(\ell\cdot n/\log{n})$ bits, by Csirmaz~\cite{Csirmaz97} (improving 
on~\cite{CapocelliSGV93}).\footnote{In a follow-up work,
Csirmaz~\cite{csirmaz1996dealer}  proved a lower bound of 
$\Omega(\ell\cdot n^2/ \log n)$ for the total share size of all parties.}

Bridging the exponential gap between
%
the upper and the lower bounds  %
is the major open
%
problem  %
in the study of secret sharing schemes. While it is widely believed
that the lower bound should be exponential 
(see, \eg, \cite{BeimelPHD,Beimel11}), no major progress has been
%
made  %
in the last two decades. Moreover, even a 
non-explicit linear lower bound is not known, that is, whether there
\emph{exists} an access structure that requires linear size shares.

\subsection{Our results}

\paragraph{Share size lower bound.} We close the gap in share size for threshold
secret sharing up to a small additive constant.  We assume for simplicity that
all parties are given equally long shares.

\begin{theorem}	\label{cor:main}
For every $n\in\mathbb{N}$ and $1 < t < n$, any $(t,n)$-threshold secret 
sharing scheme for a
1-bit secret requires shares of at least $\log (t+1)$ bits.
\end{theorem}

The assumption $1 < t < n$ is necessary, as $(1, n)$-threshold and $(n, n)$-threshold secret sharing schemes with share size 1 do exist.

Our bound is tight when $t = n - 1$ and $n$ is the power of a prime; see
\expref{Section}{sec:shamir}. By combining \expref{Theorem}{cor:main} with the lower
bound of Kilian and Nisan, we determine the share size of threshold schemes up
to a small additive constant.  That is, we get that any such scheme requires
shares of size
\begin{equation}
\label{eqn:combined}
\max\{\log(n-t+2), \log (t+1)\} \geq \log \frac{n+3}{2}.
\end{equation}

\expref{Theorem}{cor:main} is a special case of the following theorem, which applies more generally to ramp schemes. 
\begin{theorem}
	\label{thm:main}
	For every $n\in\mathbb{N}$ and $1 \leq s < r < n$, any $(s,r,n)$-ramp secret sharing scheme
	for a 1-bit secret requires shares of at least $\log ((r + 1)/(r-s))$ bits.
\end{theorem}
By combining \expref{Theorem}{thm:main} with the lower bound of~\cite{PueyoCX13}, we
get that any $(s,r,n)$-ramp secret sharing scheme must have share size at least
\begin{equation}
\label{eqn:combined_ramp}
\max\left\{\log\frac{n-s+1}{r-s}, \log \frac{r + 1}{r - s}\right\} \geq \log
\frac{n+r-s+2}{2\cdot (r-s)}.
\end{equation}

%
\paragraph{Proof technique and limitations.}
We prove our lower bounds by analyzing a new game-theoretic relaxation of secret
sharing. Here, we focus on threshold schemes, although our argument also applies to
ramp schemes.

Given an access structure $\calA$ and a real-valued parameter $\theta > 0$ we consider the following zero-sum game $G(\mathcal{A}, \theta)$:  Alice and Bob pick sets $A$ and $B$ in the access structure $\mathcal{A}$, respectively, and the payoff is $(-\theta)^{\abs{A \setminus B}}$, where $A \setminus B$ denotes set difference.  We say Alice wins if she has a strategy with non-negative expected payoff, and Bob wins otherwise.

We show (in \expref{Lemma}{lemma:main}) that if Bob wins in the game $G(\calA, 1/(q-1))$, then no secret sharing scheme with share size $\log q$ exists.  We prove \expref{Theorem}{thm:main} by constructing such a strategy for Bob.

On the negative side, we show that our analysis is optimal for threshold access structures, so the lower bound in \expref{Theorem}{cor:main} is tight with respect to this method:

\begin{theorem}	\label{thm:limitation}
For all $1 < t < n$ and $0 < \theta \leq 1/t$,  Alice wins in the game $G(\thr{t}{n}, \theta)$.
\end{theorem}

We also show that, for any total access structure $\calA$, this method cannot prove a lower bound exceeding 
$\log\abs{\min\calA}\leq \log \binom{n}{\floor{n/2}} = n-\Omega(\log{n})$, 
where $\min\calA=\{A\in\calA:\forall B\in\calA, B\not\subset A\}$ is the set of min-terms in $\calA$. 

\begin{theorem}	\label{thm:limitation1}
For every access structure $\mathcal{A}$ and every 
$0 < \theta \leq 1/(\abs{\min\calA} - 1)$ 
Alice wins in the game $G(\mathcal{A}, \theta)$. 
\end{theorem}


\subsection{Related work}
\paragraph{Known frameworks for proving lower bounds.}
The method of Csirmaz~\cite{Csirmaz97} is a 
general framework for proving lower bounds on share size in various access
structures.
Csirmaz's framework is a linear programming relaxation whose variables are the entropies of the joint distributions of the shares (one for each subset of the parties) and whose constraints are entropy inequalities.  A solution to the dual linear program yields an $\Omega(n/\log n)$ lower bound on the \emph{information ratio}, namely the length of the largest share divided by the length of the secret.

Csirmaz's framework does not give any non-trivial lower bounds on
share size for a 1-bit secret over threshold access
structures.  The reason is that for sufficiently long secrets, the information ratio of threshold schemes is 1 in Shamir's construction (see \expref{Claim}{claim:shamiropt}).  Kilian and
Nisan's~\cite{KilianN90} proof is the only known argument for threshold schemes.
It does not appear to be useful for any other access structure, including the
$(t,n)$-threshold access structures with $t$ being close to $n$.

Csirmaz~\cite{Csirmaz97} showed that his framework cannot be used to
prove a super-linear lower bound on share size for any access structure when the constraints are the Shannon inequalities from information theory. This 
claim was strengthened by Beimel and Orlov~\cite{BeimelO11}, who showed that certain additional 
``non-Shannon type'' information inequalities do not help bypass the linear share
size barrier (see~\cite{MolleviPY16} for a follow-up).

\paragraph{Linear schemes.}
A secret sharing scheme is \emph{linear} if the reconstruction procedure is a
linear function of the shares (over a finite field). Many known
schemes are linear (see~\cite{BeimelI01,BeimelW05,KomargodskiNY16,VaikuntanathanV15,BIVW,LiuVW17,LiuV18} for exceptions) and
super-polynomial lower bounds for linear schemes were given in~\cite{BabaiGW99,BeimelGP97,Gal01,pitassi-robere} via their equivalence to monotone span programs~\cite{KarchmerW93}. 

For \emph{linear} $(2,n)$-threshold secret sharing schemes for a 1-bit secret, a
$\log n$ lower bound on share size was proven by Karchmer and
Wigderson~\cite{KarchmerW93}. This was generalized by Cramer et al.~\cite{CramerF02} to a lower bound as in Equation \eqref{eqn:combined}.\footnote{Their proof is based on the duality of linear secret sharing schemes, by which the optimal share sizes of linear $(s, r, n)$-ramp schemes and linear $(n-r, n-s, n)$-ramp schemes are equal.}
For \emph{linear} $(s,r,n)$-ramp secret sharing schemes, 
Cascudo et al.~\cite{PueyoCX13} obtained a
lower bound as in Equation \eqref{eqn:combined_ramp}. We emphasize that our
lower bounds match the lower bounds of~\cite{PueyoCX13} but are not restricted
to linear (ramp) secret sharing schemes.

\section{Access structures and secret sharing}
\label{sec:prelim}

%
Let $\parties \triangleq \{1,\dots,n\}$ be a set of $n$ parties. A
collection of subsets $\calA\subseteq 2^\parties$ is \emph{upward-closed} 
if for every $B\in \calA$ and $C\supseteq B$ it holds that $C\in
\calA$.  The collection is \emph{downward-closed} if for every 
$B\in \calA$ and $C \subseteq B$ it holds that $C\in \calA$.

\begin{definition}   \label{def:accessstructure}
%
An \emph{access structure} is a pair $\calA = (\calS, \calR)$
where
\begin{itemize}
\item[(i)] $\calS$ and $\calR$ are non-empty disjoint subsets of $2^\parties$,
\item[(ii)] $\calR$ is upward-closed and $\calS$ is downward-closed.
\end{itemize}  
Subsets in $\calR$ are called \emph{qualified} and subsets
in $\calS$ are called \emph{unqualified}.
\end{definition}
\begin{remark}    \label{rem:closed}
The condition that $\calR\neq\emptyset$
is equivalent to saying that $\parties\in\calR$, \ie,
all parties together are qualified (to unlock the
secret).  The condition $\calS\neq\emptyset$ 
is equivalent to saying that $\emptyset\in\calS$, \ie,
the secret is not public (the empty set is
not qualified).
\end{remark}  

The access structure is \emph{total} if $\calR$ and $\calS$ form a partition of
$2^{\parties}$.  If $\calA = (\calS, \calR)$ is total we write $R \in \calA$ for
$R \in \calR$ and $S \not\in \calA$ for $S \in \calS$.

In this paper we are primarily concerned with 
the following two types of access structures:

\begin{itemize}
\item The \emph{threshold access structure} $\thr{t}{n}$ is a total access structure over $n$ parties in which any $t$ parties can reconstruct and secrecy is guaranteed against any subset of $t - 1$ parties:
\[ \calS = \{S\colon \abs{S} \leq t - 1\}\,, \qquad \calR = \{R\colon \abs{R} \geq t\}. \]
\item More generally, in the \emph{ramp access structure} $\ramp{s,r}{n}$, any $r$ parties can reconstruct and secrecy is guaranteed against any $s$ parties:
\[ \calS = \{S\colon \abs{S} \leq s\}\,, \qquad \calR = \{R\colon \abs{R} \geq r\}. \]
\end{itemize}

A secret sharing scheme involves a dealer who has a secret, a set of $n$
parties, and an
%
%
access structure $\calA = (\calS,\calR)$. A secret sharing scheme for $\calA=(\calS,\calR)$ is a method by
which the dealer distributes shares to the parties such that any subset in
$\calR$ can reconstruct the secret from its shares, while any subset in $\calS$
cannot infer any information on the secret.  We restrict our definition to
1-bit secrets.

\begin{definition}[Secret sharing]\label{def:SecretSharing}
A secret sharing scheme of a 1-bit secret for an 
%
%
access structure
$\calA=(\calS,\calR)$ over $n$ parties over share alphabet $\Sigma$ is
%
a pair $(p_0,p_1)$ of probability distributions %
over $\Sigma^n$ with the following properties:
\begin{description}
\item {\bf Reconstruction:} \label{def:ss_correctness} For every $R \in \calR$
the marginal distributions\footnote{The marginal distribution of $p$ over $S$ is the function $p_S\colon \mathbb{Z}_q^S \to \mathbb{R}$ given by $p_S(z) = \sum_{x\colon x_S = z} p(x)$.} 
%
%
of $p_0$ and $p_1$ on the set $R$ are disjoint.
\item {\bf Secrecy:} \label{def:ss_secrecy} For every $S \in \calS$ the marginal
distributions of $p_0$ and $p_1$ on the set $S$ are identical.
\end{description}
\end{definition}

An implementation of a secret sharing scheme consists of a sharing algorithm
that samples the shares from the probability distribution $p_0$ or $p_1$ depending on the value
of the secret and of a reconstruction algorithm that recovers the secret from
the joint values of the shares of any qualified subsets of parties. The disjointness requirement ensures that recovery by qualified subsets of parties is possible with probability 1.  The
secrecy requirement ensures that unqualified subsets of parties can extract no
information about the secret. Thus, our definition is equivalent to the ones
given, for example, in~\cite[Definition 3.6]{BeimelPHD} and in~\cite[Definitions
2 and 3]{Beimel11}.

\paragraph{An alternative formulation of secret sharing.}
Here is an equivalent formulation of secret sharing.  For $x \in \Z_q^n$,
we use $[x]$ to denote the set of positions with non-zero entries,
\ie, $[x] = \{i\colon x_i \neq 0\}$, and $\compl{[x]}$ for
the complement of this set, \ie, $\compl{[x]}=\{i\colon x_i = 0\}$.  
In this notation, $[x - y]$ is the set of positions
where  $x$ and $y$ differ and $\compl{[x - y]}$ is the set of positions
where they agree.

A function $\phi_S \colon \Z_q^n \to \mathbb{C}$ is an \emph{$S$-junta} if the value $\phi_S(x_1, \dots, x_n)$ is determined by the inputs $x_i\colon i \in S$.

\begin{lemma}
\label{lemma:sdp}
A secret sharing scheme for a 1-bit secret for
%
%
an access structure
$\calA=(\calS,\calR)$ over share alphabet $\Z_q$ exists if and only if 
there exists a function $f\colon \Z_q^n \to \R$ that is not 
identically zero satisfying the following 
%
conditions.    %
\begin{description}
\label{en:SDP}
\item \textbf{Reconstruction:} For all $x, y \in\mathbb{Z}_q^n$ such that
$\compl{[x - y]} \in \calR$,
we have  $f(x)\cdot f(y) \geq 0$.
\item \textbf{Secrecy:} For every $S \in \calS$ and every $S$-junta
$\phi_S\colon \Z_q^n \to \mathbb{C}$, we have  
$\E[f(x) \phi_S(x)] = 0$, where the
expectation is over the uniform probability distribution of $x\in\Z_q^n$.
\end{description}
\end{lemma}
\begin{proof}
For a secret sharing scheme $(p_0,p_1)$ 
we set $f(x) = p_0(x) - p_1(x)$.  The functions $p_0$ and $p_1$ have disjoint
support (otherwise even reconstruction by all parties is impossible) 
so $f$ cannot be identically zero.  The reconstruction implies that if $\compl{[x - y]} \in \calR$, then at least one of $p_0$ and $p_1$ must assign zero probability to both $x$ and $y$, so $f(x) \cdot f(y)$ equals either $p_0(x) \cdot p_0(y)$ or $(-p_1(x)) \cdot (-p_1(y))$.  In either case $f(x) \cdot f(y) \geq 0$.
For secrecy, since $p_0$ and $p_1$ have the same marginals on 
$S \in \calS$, $\E[p_0(x) \phi_S(x)]  = \E[p_1(x) \phi_S(x)]$ so 
$\E[f(x) \phi_S(x)] = 0$.
	
In the other direction, let $p_0(x) = C\cdot \max\{f(x), 0\}$ and let
$p_1(x) =C\cdot \max\{-f(x),\allowbreak 0\}$, where 
$1/C = (1/2) \sum_x \abs{f(x)}$. 
We show reconstruction by contrapositive: If $p_0$ and $p_1$ did not have
disjoint support on some set $R \in \calR$, there would exist $x, y \in
\Z_q^n$ such that $p_0(x) > 0$, $p_1(y) > 0$, and $R \subseteq \compl{[x -
y]} \in \calR$  (by the monotonicity of $\calR$), implying $f(x) > 0$,
$f(y) < 0$, and therefore $f(x) \cdot f(y) < 0$.
For secrecy, by construction we have 
$f = (p_0 - p_1)/C$, so $\E[p_0(x) \phi_S(x)] = \E[p_1(x) \phi_S(x)]$ 
for every test function $\phi_S$ that only depends on the positions in 
$S \in \calS$.  Since no $\phi_S$ can distinguish between $p_0$ and 
$p_1$ on $S$, the statistical distance between the marginal distribution 
of $p_0$ and $p_1$ on $S$ is zero, so the two are identical.
\end{proof}

By \expref{Lemma}{lemma:sdp}, ruling out the existence of a secret sharing scheme is equivalent to refuting the solvability of a real-valued constraint satisfaction problem with quadratic constraints.  
The class of refutations that we investigate are sum-of-squares inequalities of the form $\sum_i (\sum_x \rho_i(x) f(x))^2 < 0$, \ie, 
dual solutions to the canonical semidefinite relaxation of this problem.

Owing to the symmetries inherent in secret sharing, the sum-of-squares certificate of interest is a positive \emph{linear} combination of the squared magnitudes of the Fourier coefficients of $f$.  Our lower bound can be expressed as a dual certificate for a linear program with variables $\abs{\hat{f}(a)}^2, a \in \Z_q^n$.  This linear program has the form of a zero-sum game, which we describe next.


\section{A zero-sum game and proof of \expref{Theorem}{thm:main}}
\label{sec:game}

Given an access structure $\calA = (\calS, \calR)$ and a real parameter
$\theta > 0$ we define the following zero-sum game $G(\calA, \theta)$ between
Alice and Bob. The actions are a set $A \not\in \calS$ for Alice and a set
$B\in \calR$ for Bob.  
Each party chooses its action independently of the other.
The payoff of the game is $(-\theta)^{\abs{A \setminus B}}$.  We say that
Alice wins if she has a strategy with non-negative expected payoff (over her randomness) against all strategies of Bob. Analogously, we say that
Bob wins if he has a strategy with negative expected payoff.  By von Neumann's
minimax theorem the game has a unique winner. 

\begin{lemma}
	\label{lemma:main}
	If there exists a secret sharing scheme for $\calA$ with alphabet size
	$q\in\mathbb{N}$, then Alice wins in the game $G(\calA, 1/(q-1))$.
\end{lemma}

Our proof of \expref{Lemma}{lemma:main} uses Fourier analysis, which we briefly
recall here. The characters of the group $\Z_q^n$ are the complex-valued
functions $\chi_a\colon \Z_q^n \to\mathbb{C}$, where $a$ ranges over $\Z_q^n$,
defined as $\chi_a(x) = \omega^{\ip{a, x}}$, $\omega = \eee^{2\pi i/q}$.  
The
characters are an orthonormal basis with respect to the inner product $\langle
f, g \rangle = \E_{x}[f(x) \cdot \overline{g(x)}]$ with $x$
chosen uniformly from $\Z_q^n$.  The characters inherit the group structure: $\chi_a \cdot \chi_b = \chi_{a+b}$ and $\chi_a^{-1} = \overline{\chi_a} = \chi_{-a}$. Every function $f\colon \Z_q^n \to \mathbb{C}$
can then be uniquely written as a linear combination
\[
  f = \sum_{a \in \Z_q^n}
  \hat{f}(a) \cdot \chi_a
\]
with the Fourier coefficients $\hat{f}(a)$ given by
$\hat{f}(a) = \langle f,\chi_a \rangle = \E_{x}[f(x)\cdot
\overline{\chi_a(x)}]$.

\begin{proof}[Proof of \expref{Lemma}{lemma:main}]
	We show that Alice has a winning strategy. That is, we show that Alice has a strategy such that for
	every possible action of Bob, the expected payoff of the game is non-negative.  
	
	We identify the alphabet with the elements of the group $\Z_q$. Let $f\colon \Z_q^n \to \R$ be the function $f(x) = p_0(x) - p_1(x)$.  Alice plays set $A$ with probability proportional to
        \[
          \sum_{a\colon [a] = A} \abs{\hat{f}(a)}^2.
        \]
        By the secrecy part of \expref{Lemma}{lemma:sdp},
        $\E[f(x) \cdot\overline{\chi_a(x)}] = 0$ whenever
        $[a] \in \calS$, so Alice's strategy is indeed supported on
        sets outside $\calS$.
	
	Now let $B$ be an arbitrary set in $\calR$.  By the reconstruction part of
	\expref{Lemma}{lemma:sdp} and the fact that $f$ is real-valued, for every
	$x\in\Z_n^q$ and every $z\in\Z_n^q$ such that $\compl{[z]} = B$, we have that
	\begin{equation}
	\label{eq:recon}
	f(x) \cdot \overline{f(x - z)} = f(x) \cdot f(x - z) \geq 0.
	\end{equation}
	Let $x$ be uniform in $\Z_q^n$ and $z$ be uniform in $\Z_q^n$ conditioned on $\compl{[z]} = B$.  Averaging over this distribution, we have
	\begin{align*}
	\E_{x,z}[f(x) \cdot \overline{f(x - z)}] 
	&= \sum_{a, b\in\Z_q^n} \hat{f}(a) \cdot \overline{\hat{f}(b)} \cdot\E_{x, z}[\chi_a(x)
	\cdot \overline{\chi_b(x - z)}] \\
	&= \sum_a \abs{\hat{f}(a)}^2 \cdot \E_z[\chi_a(z)] \\
	&= \sum_a \abs{\hat{f}(a)}^2 \cdot \prod_{i \in [a]} \E_z[\omega^{a_iz_i}],
	\end{align*}
	where the first equality follows by writing $f(x)$ and $\overline{f(x-z)}$ using
	their Fourier representation and using linearity of expectation, the second
	equality follows since $x$ and $z$ are independent and since $\E_{x}[\chi_a(x)
	\cdot \overline{\chi_b(x)}] = 0$ for $a\neq b$, and the last equality follows
	since $z$ is chosen from a product distribution.
	
	The expression $\E[\omega^{a_iz_i}]$ evaluates to one when $i$ is in $B$ (since
	$z_i$ is fixed to zero).  Otherwise, $z_i$ is uniformly distributed over the set
	$\Z_q \setminus \{0\}$ and
	\[ \E_z[\omega^{a_iz_i}] = \frac{1}{q - 1} \sum_{z_i \in \Z_q \setminus \{0\}} \omega^{a_iz_i} = 
	\frac{1}{q - 1}\Bigl(\sum\nolimits_{z_i \in \Z_q} \omega^{a_iz_i} - 1\Bigr) = -\frac{1}{q-1}. \]
	Therefore, $\prod_{i \in [a]} \E_z[\omega^{a_iz_i}] = (-1/(q-1))^{\abs{[a] \setminus B}}$, and by Equation~\eqref{eq:recon}
	\[\sum_a \abs{\hat{f}(a)}^2 \cdot \left(\frac{-1}{q-1}\right)^{\abs{[a] \setminus B}} \geq 0.\]
Grouping all elements $a$  for which $[a] = A$, we get that 
\[ \sum_A \Bigl(\sum\nolimits_{a\colon [a] = A} \abs{\hat{f}(a)}^2\Bigr) \cdot \Bigl(-\frac{1}{q-1} \Bigr)^{\abs{A \setminus B}} \geq 0 \qquad\text{for all $B \in \calR$.} \]
Therefore, Alice's strategy has non-negative expected payoff with respect to every possible action of Bob.
\end{proof}

\begin{proof}[Proof of \expref{Theorem}{thm:main}]
	It is sufficient to prove \expref{Theorem}{thm:main} in the case $n = r + 1$:  If a secret sharing scheme for $\ramp{s, r}{n}$ existed, then a secret sharing for $\ramp{s, r}{r + 1} $ over the same alphabet can be obtained by discarding the remaining $n - r - 1$ parties and their shares.
	
	We now give a winning strategy for Bob in the game $G(\ramp{s, r}{r + 1}, \theta)$ for any $\theta > (r - s)/(s + 1)$.  By \expref{Lemma}{lemma:main} it then follows that no secret sharing scheme over an alphabet of size $(r + 1)/(r - s)$ exists.
	
	Bob's strategy is to uniformly choose a set $B$ of size $r$ (which is in $\calR$).  Then for every set $A \not\in \calS$, either $A \subseteq B$ and then $\abs{A\setminus B} = 0$, or
	$A\not\subseteq B$ and then $\abs{A\setminus B} = 1$ (since $B$ includes all
	parties except one). Thus, for every $A\not\in \calS$, the expected payoff is
	\begin{align}
	\label{eqn:analysis}
	\E_B\bigl[(-\theta)^{\abs{A \setminus B}}\bigr] 
	&= 1 \cdot \pr_B[A \subseteq B] - \theta \cdot \pr_B[A \not\subseteq B] \notag \\
	&= 1 \cdot \frac{r + 1 - \abs{A}}{r + 1} - \theta \cdot \frac{\abs{A}}{r + 1} \notag \\
	&\leq \frac{r - s}{r + 1} - \theta \cdot \frac{s + 1}{r + 1},
	\end{align}
	where the inequality follows since $\abs{A} \geq s+1$.  If $\theta > (r - s)/(s + 1)$ this expression is less than zero, \ie, Bob wins.
\end{proof}

It is also possible to deduce \expref{Theorem}{thm:main} directly from
\expref{Lemma}{lemma:main} by showing the existence of a winning strategy for Bob in
the game $G(\ramp{s, r}{n}, \theta)$ whenever $\theta > (r - s)/(s + 1)$ (rather
than for $G(\ramp{s, r}{r+1}, \theta)$, as we did above).  Let $R$ be a random subset of $r + 1$ parties.  Bob's strategy has the form $B = B_0 \cup B_1$, where $B_0$ is a uniformly random subset of $R$ of size $r$ and $B_1$ is a random subset of $\compl{R}$ obtained by including each element independently with probability $p = \theta/(1 + \theta)$.  The value of $p$ is chosen so that a random variable that equals $1$ with probability $p$ and $-\theta$ with probability $1 - p$ is unbiased.

Let $A$, where $\abs{A} \geq s + 1$, be any action of Alice.  For a fixed choice
of $R$, if $A \setminus R$ is nonempty, by our choice of probability $p$ the
expected payoff is zero.  Otherwise, $A$ is a subset of $R$, and by Equation \eqref{eqn:analysis} the expected payoff is at most $-(s + 1) \cdot \theta + (r - s) < 0$.  Since the event $A \subseteq R$ has positive probability, the expected payoff is negative and Bob wins.

\paragraph{Extension to unequal shares.}
\expref{Theorem}{cor:main} requires that the shares given to all parties have the same length. Its proof extends easily to yield the following
generalization.  %
For every $n$, every $1 < t < n$, and every $(t, n)$-threshold secret sharing scheme in which party $i$ receives a $\log q_i$-bit share and $q_1 \leq q_2 \leq \dots \leq q_n$ it must hold that
\begin{equation}
\label{eqn:unequal}
\frac{1}{q_1} + \dots + \frac{1}{q_{t+1}} \leq 1.
\end{equation}
In particular, inequality \eqref{eqn:unequal} implies that the \emph{average}
share size must be at least $\log{(t+1)}$.   Kilian and Nisan~\cite{KilianN90}
prove the same for $(n - t + 1, n)$-threshold access structures.

The proof of inequality~\eqref{eqn:unequal} is a direct extension of the proof of \expref{Theorem}{cor:main}.     We describe the differences and main steps.
Given an access structure $\calA = (\calS, \calR)$ and
real positive numbers $\theta_1,\dots, \theta_n$, we define the following zero-sum game $G(\calA, \theta_1,\dots,\theta_n)$ between Alice and Bob. The actions are a set $A \not\in \calS$ for Alice and a set $B
\in \calR$ for Bob.  The payoff of the game is $\prod_{i \in A \setminus B} (-\theta_i)$.  
We can obtain the following generalized version of \expref{Lemma}{lemma:main}.

\begin{lemma}
\label{lemma:main_ave}
If there exists a secret sharing scheme for $\calA$ with alphabet size $q_i\in\mathbb{N}$ for the $i$-th party,  then Alice wins in the game $G(\calA, 1/(q_1-1),\dots,1/(q_n-1))$.
\end{lemma}

The generalized lemma can be proved via Fourier analysis over the product group $\Z_{q_1} \times \cdots \times \Z_{q_n}$.    We omit this proof.

As in the proof of \expref{Theorem}{cor:main}, it is sufficient to establish inequality~\eqref{eqn:unequal} in the special case $t = n - 1$.  In this case, Bob plays set $B = \{1, \dots, n\}\setminus\{i\}$ with probability proportional to $1 - 1/q_i$.    It can be verified that when $\sum_{i=1}^n 1/q_i > 1$ this is a winning strategy for Bob.  
%
Specifically,   for every set $A \not\in \calS$, either $A \subseteq B$ and then $\abs{A\setminus B} = 0$, or $A\not\subseteq B$ and then $\abs{A\setminus B} = 1$ (since $B$ includes all parties except one).  Let $C=\sum_{i=1}^n (1 - 1/q_i)$.  Thus, for every $A\not\in \calS$, the expected payoff is 
\begin{align*}
\label{eqn:analysis_ave}
\E_B\bigl[\prod_{i \in A \setminus B} -\theta_i\bigr] 
&=  \pr_B[A \subseteq B] \cdot 1- \sum_{i\in A} \pr_B[A \setminus B = \{i\}]\cdot (-\theta_i) \notag \\
&= \sum_{i\notin A}\frac{1-1/q_i}{C}  \cdot  1- \sum_{i\in A}\frac{1-1/q_i}{C} \cdot \theta_i \notag \\
& =\sum_{i\notin A}\frac{1-1/q_i}{C}- \sum_{i\in A}\frac{1/q_i}{C} \\
&= 1 - \frac{|A|}{C}.
\end{align*}
By assumption $\abs{A} \geq n-1$. So if $\sum_{i=1}^n 1/q_i > 1$ this expression is less than zero, \ie, Bob wins.


\section{Limitations of the game relaxation}
\label{sec:limitation}

In the case of threshold access structures \expref{Theorem}{thm:main} shows that Bob
has a winning strategy in the game $G(\thr{t}{n}, \theta)$ whenever $\theta >
1/t$.  We now prove \expref{Theorem}{thm:limitation}, which states that our analysis
is optimal:  There exists a winning strategy for Alice when $\theta \leq 1/t$.  

We also prove \expref{Theorem}{thm:limitation1}: For every total access structure $\calA$ over $n$ parties, Alice has a winning strategy in $G(\calA, \theta)$ for every $\theta \leq 1/(\abs{\calA}-1)$.  As the proof of \expref{Theorem}{thm:limitation1} is simpler, we present that one first. We remark that 
\expref{Theorem}{thm:limitation1} can be generalized to any
access structure $(\calS,\calR)$ by replacing $\calA$ with $\calR$ in the proof.


\begin{proof}[Proof of \expref{Theorem}{thm:limitation1}]
Alice's strategy is uniformly random over all minterms $A \in \min\calA$.  
Then, for every $B \in \calA$ and $\theta<1$, it holds that
  \begin{align*}
  \E_A[(-\theta)^{\abs{A \setminus B}}] &= \E_A[(-\theta)^{\abs{A \setminus B}} \mid A\subseteq B]\cdot \pr_A[A\subseteq B] \\ &\qquad + \E_A[(-\theta)^{\abs{A \setminus B}} \mid A\not\subseteq B]\cdot \pr_A[A\not\subseteq B] \\
	& \geq 1 \cdot \pr_A[A \subseteq B] - \theta \cdot
	\pr_A[A \not\subseteq B]\\ 
	& = (1+\theta)\cdot \pr_A[A \subseteq B] - \theta \\
	& \geq (1+\theta)\cdot \frac{1}{\abs{\min\calA}} - \theta.
	\end{align*}
This is non-negative when $\theta \leq 1/(\abs{\min\calA} - 1)$.
\end{proof}

\begin{proof}[Proof of \expref{Theorem}{thm:limitation}]
Let $a_0, \dots, a_n$ be the following sequence of integers:
\[ a_0=\dots=a_{t-1}=0, \quad a_t=1, \quad a_s = k_t \cdot a_{s-1} + \dots + k_0 \cdot a_{s-t-1} \]   
for $t+1\leq s\leq n$, where $k_j$ is the coefficient of $x^j$ in the formal expansion of $(x+1)^t\cdot
(1/\theta - x)$.
By expanding this expression according to the Binomial formula, we see
that the numbers $k_0, \dots, k_t$ are non-negative when $\theta \leq 1/t$ because
\[ k_j = \binom{t}{j}\Bigl(\frac{1}{\theta} - \frac{j}{t-j+1}\Bigr) \geq 0 \]
for all $0\leq j\leq t$.  Therefore $a_s$ is also non-negative for all $s$.
	
Alice plays set $A$ with probability proportional to the number $a_{\abs{A}}$.  We will prove that this is a winning strategy for Alice.  When $B = \{1, \dots, n\}$, then $\E_A[(-\theta)^{\abs{A \setminus B}}] = 1$ and Alice wins.  Now let $B\subseteq \{1, \dots, n\}$ be any set such that $t \leq \abs{B} < n$. Let
	\[ \theta_j = \begin{cases}1, &\text{if $j \in B$}, \\ -\theta, &\text{if $j \not\in B$}.\end{cases} \]
	%
	Then, $\E_A[(-\theta)^{\abs{A \setminus B}}]$ is proportional to 
	\[  \sum_A a_{\abs{A}} \prod_{j \in A} \theta_j = \sum_{s = 0}^n a_s w_s \quad\text{where}\quad w_s = \sum_{A\colon \abs{A} = s} \prod\nolimits_{j \in A} \theta_j. \]
The number $w_s$ can be represented as the coefficient of $z^s$ in the formal expansion of $g_0(z) = \prod_{j=1}^n (1 + \theta_j z)$.  Since exactly
$\abs{B}$ of the $\theta_j$  
equal $1$ and the other $n - \abs{B}$ equal $-\theta$, it follows that
	\begin{equation} \label{eq:g_t_closed_form} g_0(z)=(1+z)^{\abs{B}}
	\cdot(1 - \theta z)^{n-\abs{B}}.
	\end{equation}
	The numbers $a_0,\dots,a_n$ (as defined in the beginning of the proof) are defined by an order-$t$ homogeneous linear
%
recurrence %
with constant coefficients whose characteristic equation is
$(x + 1)^t \cdot (1/\theta - x) = 0$. 
This equation has roots $-1$ (with multiplicity $t$) and $1/\theta$ 
(with multiplicity $1$). Therefore,
\[ a_s = C \cdot \theta^{-s} + \sum_{i=0}^{t-1} c_i \cdot s^i \cdot (-1)^s \]
where $c_0, \dots, c_{t-1}$ and $C$ are constants determined by the initial conditions on $a_0, \dots, a_t$.  We can now write
\[ \sum_{s=0}^n a_s \cdot w_s = C \cdot \sum_{s=0}^n w_s \cdot \theta^{-s} +
\sum_{i=0}^{t-1} c_i \cdot \sum_{s=0}^n w_s \cdot s^i \cdot (-1)^s. \] As
$g_0$ is the generating function of 
the  %
$w_s$, $g_0(z) =
\sum^n_{s=0} w_s \cdot z^s$. So, the term $\sum_{s=0}^n w_s \cdot \theta^{-s}$
equals $g_0(1/\theta) = 0$.  To finish the proof, we show that $\sum_{s=0}^n w_s
	\cdot s^i \cdot (-1)^s = 0$ for all $i \leq t - 1$. This implies that Alice's
	strategy has a 0 payoff, so she wins the game.  Let $g_i(z) =
	z\cdot g'_{i-1}(z)$ for $1 \leq i \leq t - 1$ where $g'_{i-1}$ is the derivative
	of $g_{i-1}$.  On the one hand, since $-1$ is a root of $g_0$ of multiplicity
	$t$, $g_i(-1) = 0$ for all $i \leq t - 1$.  On the other hand, $g_i(z)$ has the
	formal expansion $\sum_{s=0}^n w_s \cdot s^i \cdot z^s$.  Therefore,
	$\sum_{s=0}^n w_s \cdot s^i \cdot (-1)^s$ must equal zero.
\end{proof}

\section{On the tightness of \expref{Theorem}{thm:main}}
\label{sec:shamir}

We show that \expref{Theorem}{thm:main} is tight when $t = n - 1$ and $n$ is
a prime power. This result is known (see, \eg, \cite[Theorem 11.13]{CDN2015})
and we give it here for completeness.

\begin{claim}
	\label{claim:shamiropt}
For every prime power $n$ 
there exists a $(n-1)$-out-of-$n$ secret sharing scheme for 
$\lfloor \log n\rfloor $-bit secrets with $\lceil \log n\rceil $-bit shares.
\end{claim}

\expref{Claim}{claim:shamiropt} follows by an optimization 
of Shamir's secret sharing scheme, where we place the secret as
the coefficient of the highest-degree term (instead of as the constant term).  We give the construction and sketch the 
proof of correctness.

To share a secret $s \in \F_n$, let $p(x) = sx^{n-2} + r(x)$, where $r$ 
is a random polynomial of degree $n - 3$ and all algebra is over the 
finite field $\F_n$.  The shares are the $n$ values $p(x)$ as $x$ 
ranges over $\F_n$. Reconstruction is immediate as the polynomial 
$p$ can be interpolated from any $n - 1$ of its values.

%
For secrecy, we show for any $s \in \F_n$ and distinct $x_1,\dots,x_{n-2} \in \F_n$, the vector $(p(x_1), \dots, p(x_{n-2}))$ is uniformly random in $\F_n^{n-2}$.  Since $p(x) = sx^{n-2} + r(x)$ it suffices to show that $(r(x_1),\dots, r(x_{n-2}))$ is uniformly random.  This is true because the evaluation map that takes the coefficients of $r$ to  
its values $r(x_1),\dots, r(x_{n-2})$ is a full-rank Vandermonde matrix.

\section{Concluding remarks}

\paragraph{Relation to the lower bound of Kilian and Nisan.}
By \expref{Theorem}{thm:limitation} our analysis of threshold secret sharing is tight within the game-theoretic relaxation that we introduce here.  As the lower bound of Kilian and Nisan~\cite{KilianN90} is incomparable with ours, their analysis cannot be cast in terms of a winning strategy in our game.  The analysis of Kilian and Nisan relies on the collision probabilities $\pr[X_A = Y_A]$, where $X$ and $Y$ are independent samples of the distributions $p_0$ and $p_1$, respectively, and $X_A, Y_A$ are their projections on $A$.  In the Fourier basis these probabilities can be written as
\[
\pr[X_A = Y_A] = \sum_{a\colon [a] \subseteq A} \hat{p}_0(a) \cdot \overline{\hat{p}_1(a)}.
\]
Our game-theoretic formulation cannot express such quantities.  It is, however, 
 possible to unify our analysis with that of Kilian and Nisan by a common linear program.  One approach is to refine the canonical linear programming formulation of the zero-sum game $G(\mathcal{A}, \theta)$, which has a variable $v(A)$ representing $\sum_{[a] = A} \abs{\hat{f}(A)}^2$ for every subset $A$ of $[n]$.  In the refined formulation, each $v(A)$ is split into $v_{00}(A) + 2v_{01}(A) + v_{11}(A)$, where $v_{ij}(A)$ represents the (real-valued) expression $\sum_{[a] = A} \hat{p}_i(a) \cdot \overline{\hat{p}_j(a)}$.  These variables satisfy $v_{ii}(A) \geq 0$ (non-negativity), $v_{00}(A) = v_{01}(A) = v_{11}(A)$ for $A \in \mathcal{S}$ (secrecy), and
\[  \sum_A v_{ij}(A) (-\theta)^{\abs{A \setminus B}} \ 
\begin{cases}
= 0, &\text{if $B \in \mathcal{R}$ and $ij = 01$,} \\
\geq 0, &\text{otherwise}.
\end{cases} 
\quad\text{(reconstruction)}
\]
The last quantity represents the probability of the event $\{b\colon X_b = Y_b\} = B$ up to scaling, where $X$ and $Y$ are independent samples of $p_i$ and $p_j$, respectively.

This linear program captures both our proof technique and that of Kilian and Nisan and is therefore infeasible with respect to $\thr{t}{n}$ when $q < \max\{n - t + 2, t + 1\}$. Our computer experiments show that this bound is tight for all $n \leq 7$, so this particular formulation does not appear to be advantageous for further improving the share size lower bounds on $\thr{t}{n}$ obtained in this paper.

\paragraph{Comparison with Csirmaz's method.}
A common feature of Csirmaz's method and ours is that the secrecy and reconstruction axioms are formulated as linear constraints on functionals of random variables determined by subsets of the shares.  Csirmaz's relevant functional is the Shannon entropy of the corresponding subset.  In contrast, under a suitable change of basis, our functionals capture the collision (R{\'e}nyi) entropies of shares of 0 and/or 1 restricted to the relevant subset of the parties.  

Csirmaz's method more generally gives a lower bound on  
the \emph{information ratio} of schemes with arbitrary secret alphabet, while our formulation only applies to binary secrets.  In this setting the two methods are incomparable.  On the one hand, our main result shows that Csirmaz's method is less powerful than ours for threshold access structures.  On the other hand, the alphabet size lower bound of 4 for the tree access structure $\{01, 12, 03, 34, 05\}$ obtained by Csirmaz and Tardos~\cite{csirmaz-tardos} cannot be matched by our game analysis.


\paragraph{On general access structures.}
We do not know what is the best possible lower bound on share size that our method can give among all access structures on $n$ parties.  \expref{Theorem}{cor:main} shows that a lower bound of $\log (n - 1)$ is attainable, while \expref{Theorem}{thm:limitation1} shows that a lower bound of $\log\binom{n}{\floor{n/2}}$ 
cannot be proved by our method. 

The best possible bound is the logarithm of
\[ b_n = \min_\text{$\calA$} \max \bigl\{q\colon \text{Bob wins in $G(\calA, 1/(q-1))$}\bigr\}, \]
where the minimum is taken over all access structures $\calA$ on $n$ parties.  As indirect evidence that $\log\binom{n}{\floor{n/2}}$  might be far from tight, we can prove that if the payoff function is replaced by $(-\theta)^{\abs{A \triangle B}}$, where $\triangle$ is symmetric set difference, then the quantity analogous to $b_n$ is at most 
$cn + 1$, where $c = 1/\ln(1 + \sqrt{2}) \approx 1.135$, 
when Bob plays tit-for-tat.

\begin{claim}
If $A, B \subseteq [n]$ are independent and identically distributed then $\E[(-1/cn)^{\abs{A \triangle B}}] > 0$.  
\end{claim}
\begin{proof}
We first show that $\pr[\abs{A \triangle B} = s] \leq \binom{n}{s} \pr[A = B]$.  Let $p$ be the probability mass function of the sets.
\begin{multline*}
\pr[\abs{A \triangle B} = s] = \sum_{A, B\colon \abs{A \triangle B} = s} p(A) p(B) 
   = \sum_{A, S\colon \abs{S} = s} p(A) p(A \triangle S)  \\ \leq \sqrt{\sum_{A, S\colon \abs{S} = s} p(A)^2} \cdot \sqrt{\sum_{A, S\colon \abs{S} = s} p(A \triangle S)^2} = \binom{n}{s} \sum_A p(A)^2,
\end{multline*}
where the inequality is Cauchy-Schwarz.  The expectation of interest is then lower bounded by

\begin{align*}
\E\bigl[(-1/cn)^{\abs{A \triangle B}}\bigr] 
   &\geq \pr[A = B] - \sum_{\text{$s$ odd}} \Bigl(\frac{1}{cn}\Bigr)^s \pr[\abs{A \triangle B} = s] \\
   &\geq \pr[A = B] - \sum_{\text{$s$ odd}} \Bigl(\frac{1}{cn}\Bigr)^s \binom{n}{s} \pr[A = B] \\
   &> \pr[A = B] \cdot \biggl(1 - \frac{c^{-1}}{1!} - \frac{c^{-3}}{3!} - \frac{c^{-5}}{5!} - \cdots\biggr) \\
   &= \pr[A = B] \cdot \biggl(1 - \frac{\eee^{1/c} - \eee^{-1/c}}{2}\biggr), 
\end{align*}
which is nonnegative by our choice of $c$.
\end{proof}

%

\subsubsection*{Acknowledgments.}
We thank Moni Naor for telling us about the work of Kilian and Nisan. We thank the anonymous reviewers for their useful advice.


\bibliographystyle{tocplain}   %

%
%
%
%
\bibliography{v016a002}

%
%
%
%
%
\begin{tocauthors}
\begin{tocinfo}[Bogdanov]
  Andrej Bogdanov\\
  Associate professor\\
  Department of Computer Science and Engineering\\
  The Chinese University of Hong Kong\\
  Hong Kong\\
  andrejb\tocat{}cse\tocdot{}cuhk\tocdot{}edu\tocdot{}hk \\ %
  \url{http://www.cse.cuhk.edu.hk/~andrejb/} %
\end{tocinfo}
\begin{tocinfo}[Guo]
  Siyao Guo\\
  Assistant professor\\
  Department of Computer Science\\
  NYU Shanghai\\
  China\\
  siyao\tocdot{}guo\tocat{}nyu\tocdot{}edu \\
  \url{https://sites.google.com/site/siyaoguo/}
\end{tocinfo}
\begin{tocinfo}[Komargodski]
  Ilan Komargodski\\
  Assistant professor and Scientist\\
  Hebrew University of Jerusalem and NTT Research\\
  Jerusalem, Israel and Palo Alto, CA, USA \\
  ilan.komargodski\tocat{}mail\tocdot{}huji\tocdot{}ac\tocdot{}il \\
  \url{http://ilan.tech.cornell.edu}
\end{tocinfo}
\end{tocauthors}

%
%
%
%
%
%
%
%
%
%
%
%
%
%
\begin{tocaboutauthors}
\begin{tocabout}[Bogdanov]  %
  \textsc{Andrej Bogdanov} is an associate professor of Computer Science at the Chinese University of Hong Kong. He is on short leave from the \href{http://www.iyengaryogahongkong.com}{Iyengar Yoga Institute of Hong Kong} after a disastrous attempt at \href{https://www.yogapedia.com/definition/8043/utthita-pada-sirsasana}{utthita pada sirsasana}.  He was advised by \href{https://lucatrevisan.wordpress.com}{Luca Trevisan} who taught him the fine art of enjoying \href{https://en.wikipedia.org/wiki/Xiaolongbao}{xiaolongbao}.
\end{tocabout}
\begin{tocabout}[Guo]
%
  \textsc{Siyao Guo} is an assistant professor of Computer Science at \href{https://shanghai.nyu.edu/}{NYU Shanghai}. She received her \phd\
from the \href{http://www.cuhk.edu.hk/english/index.html}{Chinese University of Hong Kong}, advised by 
\href{http://www.cse.cuhk.edu.hk/~andrejb/}{Andrej Bogdanov}, who taught her how to eat \href{https://en.wikipedia.org/wiki/Xiaolongbao}{xiaolongbao} properly.
\end{tocabout}
\begin{tocabout}[Komargodski]
  \textsc{Ilan Komargodski} is an assistant professor of Computer Science at the
  \href{https://www.cs.huji.ac.il/}{Hebrew University of Jerusalem} and a Scientist at \href{https://ntt-research.com/cis/}{NTT Research}. During his \phd\
at the \href{http://www.weizmann.ac.il/pages/}{Weizmann Institute of Science}, advised by
  \href{http://www.wisdom.weizmann.ac.il/~naor/}{Moni Naor}, he visited  \href{http://www.cse.cuhk.edu.hk/~andrejb/}{Andrej Bogdanov} and \href{https://sites.google.com/site/siyaoguo/}{Siyao Guo} who introduced to him  \href{https://en.wikipedia.org/wiki/Xiaolongbao}{xiaolongbao} and other delicious treats like Fourier analysis. 

%
%
%
%
%
%
%
%
%
%
%
%
%
%
%
%
%
%
\end{tocabout}
\end{tocaboutauthors}

\end{document}
